\documentclass[11pt]{article}

\usepackage{amsmath,amssymb}
\usepackage{booktabs}
\usepackage{geometry}
\geometry{margin=1.1in}

\title{From the Calibrated Model to the Estimated Model:\\
Modifications Required Beyond Subsection 4.2}
\author{}
\date{\today}

\begin{document}
\maketitle

\begin{abstract}
This note catalogs every modification made to the theoretical model in
\texttt{ShotgunPCestCovid.tex} -- as fully specified through Section~1
(Introduction), Section~2 (Budget Constraints), Section~3 (The Model),
and Subsections~4.1--4.2 (Model Summary and Calibration) -- in order to
arrive at the model actually taken to the data,
\texttt{mmt8bOccBin3\_covid\_bT\_baa.mod}, and its Bayesian estimation
results (current round: reserve-equation fix, diffuse $b_\pi,b_y$
priors, \texttt{hessian} jumping covariance, $\bar\tau=0.14$, $\chi=0.5$;
see \texttt{mmt8bOccBin3\_covid\_bT\_baa\_epsL\_summary.tex} for the full
methodology and results, including the superseded tight-prior ``gen1''
round retained there for comparison). It is intended as a working
reference for rewriting the paper's Subsection~4.3 onward, which
currently continues with the calibrated-model determinacy analysis and
an older, no-longer-accurate Section~5 estimation draft (the latter
lacks the debt-to-GDP observable, the Baa-based loan rate, the COVID
labor-supply shock, a genuine OccBin ZLB likelihood, and uses a
different fiscal calibration than the results actually obtained). No
existing files are modified by this note; it is a planning document only.
\end{abstract}

\tableofcontents

\section{Scope and a frequency caveat}

Sections~1--3 and Subsections~4.1--4.2 present the model at a
calibration that does not specify an explicit sampling frequency, with
illustrative annual-scale targets ($\Pi=1.04$, $r=1.07$, $\beta=0.972$,
etc.). The estimated model is explicitly \emph{quarterly}, disciplined
by a 2009Q1--2025Q4 U.S.\ sample, and every rate-like calibrated value
is recalibrated accordingly -- the same symbol (e.g.\ $\beta$) refers to
a materially different numerical target between the two. This is not
itemized separately below except where the reinterpretation also
changes the model's economic content (e.g.\ the ZLB), but should be
kept in mind throughout: a side-by-side value comparison for shared
symbols is only meaningful once both are understood as quarterly.

\section{Structural equation modifications}

\subsection{Interest rate rule: output response, steady-state rate base, shock, and ZLB}
Paper (eq.~\ref{ior}-equivalent, Section~2):
\[
r^{n}_{t}=r_{t}^{\alpha_{i}}\left(\frac{\Pi_{t}}{\Pi}\right)^{b_{\pi}}.
\]
Estimated model:
\[
r^{n}_{t}=\max\!\left[1,\;\big(r^{\ell,ss}\big)^{a_i}\left(\frac{\pi_t}{\pi^{ss}}\right)^{b_\pi}\left(\frac{y_t}{y^{ss}}\right)^{b_y}\exp(z^r_t)\right],
\qquad z^r_t = b_r\,z^r_{t-1}+\varepsilon^r_t.
\]
Three changes: (i) an output-response term $(y_t/y^{ss})^{b_y}$ is added
(no output term exists in the paper's rule at all); (ii) the base rate
is the \emph{fixed calibrated steady-state} loan rate $r^{\ell,ss}$
rather than the time-varying market rate $r_t$ (the paper's $r_t$ and
loan rate $r^\ell_t$ coincide via the arbitrage condition
$r^\ell_t=r_t$, but the estimated rule uses the calibrated anchor, not
the contemporaneous value); (iii) a persistent monetary-policy shock
$z^r_t$ is added, entirely absent from the paper. Most substantively,
(iv) a strict zero lower bound is imposed via a $\max[1,\cdot]$
operator, solved with Dynare's OccBin piecewise-linear toolkit
(Guerrieri and Iacoviello, 2015) -- the paper's calibrated model has no
lower bound on $r^n_t$ at all. Enforcing the ZLB genuinely (rather than
carrying the constraint syntactically but disabling it in estimation,
as an earlier draft of this file did) requires OccBin's true
piecewise-linear likelihood, which is itself nonsmooth and required a
two-stage mode-finding procedure (Section~6 below).

\subsection{Tax rule: lagged debt, persistence, shock}
\label{sec:taxrule}
Paper (eq.~\ref{tau}):
\[
\tau_t = \bar\tau\left(\frac{b_t}{(1-\Omega)b_0^T}\right)^{\varphi_b}.
\]
Estimated model:
\[
\tau_t = \tau_{t-1}^{\rho_\tau}\,\bar\tau^{\,1-\rho_\tau}\left(\frac{b_{t-1}}{(1-\omega)b_0}\right)^{\varphi_b}\exp(\varepsilon^\tau_t).
\]
Three changes: the debt ratio is \emph{lagged} ($b_{t-1}$, not $b_t$)
-- required so that $b_t$ is a predetermined state variable, which
matters for Blanchard-Kahn compatibility once $b_t$ is genuinely
endogenous (see Section~2.3); an AR(1)-type persistence term
$\tau_{t-1}^{\rho_\tau}$ is added (the paper's rule is a purely static
function of the debt ratio); and a stochastic innovation
$\varepsilon^\tau_t$ is added (the paper's rule is deterministic).

\subsection{Debt-issuance equation: linearization and simplification}
\label{sec:debtissuance}
The paper actually contains \emph{two} separate real budget-constraint
identities (Section~3, ``Equilibrium Household Budget Constraints''),
not one: eq.~(\ref{HH BC m}), which solves for currency,
$m_t = g_t - \tau_t + r^n_t n_t$, and eq.~(\ref{HH BC n}), which solves
for reserves, $n_t = m_{t-1}/\Pi_t + b_{t-1}/\Pi_t - b_t/r_t$ -- the
latter is the one \emph{including} the reserve stock $n_t$ and the
interest-on-reserves rate $r^n_t$ (via $m_t$'s own definition) explicitly,
and is discussed separately in Subsection~\ref{sec:reserve} below,
since the estimated model modifies it too, in the opposite direction
(restoring rather than dropping terms). The two identities combine into
one consolidated government budget constraint. The Policy subsection
then writes a \emph{third}, separate equation -- eq.~(\ref{b issue})
below -- that rearranges this same consolidated identity to solve
directly for $b_t$ instead of $n_t$; being a rearrangement of the same
underlying identity, it too contains $m_{t-1}$ and $n_t/r^n_t$ terms.
This subsection concerns only that third equation, the one solving for
$b_t$.

Paper (eq.~\ref{b issue}):
\[
b_t = r_t\left[b_{t-1}\Pi_t^{-1} + g - \tau_t - m_{t-1}(1-\Pi_t) + (r^n_t-1)n_t\right].
\]
Estimated model:
\[
b_t = \frac{r^{\ell,ss}}{\pi^{ss}}b_{t-1} + g_t - \tau_t + b_0\left(1-\frac{r^{\ell,ss}}{\pi^{ss}}\right) - (g-\bar\tau).
\]
This is the single largest structural departure from the paper. The
time-varying ratio $r_t/\Pi_t$ multiplying $b_{t-1}$ is replaced by a
\emph{fixed} calibrated ratio $r^{\ell,ss}/\pi^{ss}$ (an extension of
Assumption~2's fixed-real-rate logic to the dynamic debt equation); the
currency-growth term $-m_{t-1}(1-\Pi_t)$ and the reserve-interest term
$(r^n_t-1)n_t$ are dropped entirely; and government spending becomes
the stochastic $g_t$ rather than the fixed $g$. This backward-looking
approximation decouples $b_t$ from the model's forward-looking
variables ($\pi_t$, $r^\ell_t$) and is what makes a genuinely
endogenous, estimable debt stock compatible with the Blanchard-Kahn
order/rank conditions -- the full nonlinear equation (\ref{b issue})
introduces exactly the kind of forward-looking-variable coupling that a
naive endogenous-debt extension of this model class runs into. Contrast
this with the $n_t$ equation (Subsection~\ref{sec:reserve}): there, the
analogous $m_{t-1}$ and $b_{t-1}/b_t$ terms are \emph{not} dropped, only
the fixed-ratio linearization is applied where needed -- the difference
is that $b_t$'s forward-looking coupling comes from $r_t/\Pi_t$
\emph{multiplying the whole bracket} (including $g_t,\tau_t$), which
this equation's own recursive structure cannot avoid, whereas $n_t$'s
equation only ever multiplies $b_{t-1}$ and $b_t$ by already-fixed
rates, so no comparable simplification is forced on it.

\subsection{Euler equation: marginal-utility reformulation (solution technique, not economics)}
Paper (eq.~\ref{Euler}):
\[
\beta E_t\!\left(\frac{r_{t+1}}{\Pi_{t+1}}\left(\frac{h_{t+1}}{h_t}\right)^\psi\frac{w_t}{w_{t+1}}\frac{A_{t+1}}{A_t}\right)=1.
\]
Estimated model introduces $\mu_t \equiv A_t w_t h_t^\psi$ and rewrites:
\[
\mu_t = A_t w_t h_t^{\psi}, \qquad
\beta\,\frac{r^\ell_{t+1}}{\pi_{t+1}}\,\frac{\mu_{t+1}}{\mu_t}=1.
\]
Algebraically equivalent to the paper's Euler equation, not a
behavioral change. The reformulation reduces the count of
forward-looking variables in the linearized system (from what the
paper's direct specification would imply down to four: $\pi$, $r^\ell$,
$\mu$, $y$), which is what makes the model numerically tractable for
Blanchard-Kahn/E-stability analysis with an endogenous bond stock.

\subsection{Labor-leisure condition: productivity in place of the wage, plus the COVID shock}
Paper (Section~4.1, combining eqs.~\ref{foc c} and \ref{foc h}):
\[
\chi h^s_t + \frac{c_t}{w_t}h_t^\psi = \eta^{-1}.
\]
Estimated model:
\[
\chi\, hs_t + \frac{c_t}{z}h_t^{\psi} = \frac{1}{\eta}\exp(\varepsilon^L_t).
\]
Two changes: the denominator switches from the (endogenous) real wage
$w_t$ to the fixed productivity parameter $z$; and a new iid shock
$\varepsilon^L_t\sim N(0,\sigma_L^2)$ multiplies the right-hand side. A
negative realization of $\varepsilon^L_t$ shrinks the right-hand side,
forcing lower labor supply for given consumption -- a negative
labor-supply shock, included as a simple stand-in for the COVID-19
labor-supply collapse. (A deterministic one-period dummy,
\texttt{varexo\_det}, was considered first to target a specific
calendar quarter, but Dynare 7.1's estimation routine does not read
\texttt{varexo\_det} series from the data file, so it would be silently
ignored by the Kalman filter; the iid shock avoids this.) This shock
has no counterpart anywhere in the paper.

\subsection{New variable: debt-to-GDP link}
Not present in the paper at all. A purely definitional auxiliary
variable is added to connect the model's debt-to-output state $b_t$ to
an observable debt-to-GDP ratio:
\[
bT_t = \frac{b_t}{1+g}.
\]

\subsection{Reserve equation: restored debt-state coupling (update)}
\label{sec:reserve}
Unlike every other modification catalogued here, this one \emph{reduces}
the gap to the paper rather than adding to it. The estimated model's
equation for reserves $n_t$ -- the analog of the paper's
eq.~(\ref{HH BC n}), $n_t = m_{t-1}/\Pi_t + b_{t-1}/\Pi_t - b_t/r_t$ --
had been implemented as
\[
n_t = \frac{m_{t-1}}{\pi_t} + b_0(1-\omega)\left(\frac{1}{\pi_t}-\frac{1}{r^\ell_t}\right),
\]
using the fixed steady-state debt target $b_0(1-\omega)$ in place of the
actual debt state, so $n_t$ never responded to fluctuations in
outstanding debt at all. This is now corrected to
\[
n_t = \frac{m_{t-1}}{\pi_t} + \frac{b_{t-1}}{\pi_t} - \frac{b_t}{r^\ell_t},
\]
the model's direct, literal analog of eq.~(\ref{HH BC n}), with no
further simplification needed: $b_t$'s own law of motion (the fixed-ratio
linearization in Subsection~\ref{sec:debtissuance}) never depends on any forward-looking
variable, so it remains safe to use here in place of the fixed
placeholder without reopening the Blanchard-Kahn problem that motivated
that linearization in the first place. Verified directly: the steady
state and the Blanchard-Kahn eigenvalue count (4 forward-looking
variables) are unchanged before and after this fix.

\section{New stochastic structure}

\textbf{Update: the paper's Section~3 now already includes most of
this.} At the time this catalog was first written, the paper's
calibrated model (Sections~1--3, 4.1--4.2) had exactly \emph{one}
stochastic shock, the preference shock $a_t$ in the utility function.
The paper's Section~3 has since been edited (by instruction, separately
from the estimated-model changes tracked elsewhere in this note) so
that four more of these processes -- $z_t$, $v_{a,t}$, $x_{a,t}$, and
$g_t$ -- are now introduced directly in the paper itself, each at the
point where the corresponding parameter first appears (productivity
$z_t$ in the Firm decision subsection, eq.~\ref{z}; shopping-time
preference $v_{a,t}$ in the Household Problem subsection, eq.~\ref{va};
deposit technology $x_{a,t}$ in the Banking Sector subsection,
eq.~\ref{xa}; government spending $g_t$ in the Policy subsection,
eq.~\ref{g}), using the identical AR(1) functional forms as the
estimated model. \emph{Only the three rows marked ``still new'' below
remain absent from the paper}: the monetary-policy shock $z^r_t$ (which
would belong in Section~2's interest-rate rule, eq.~\ref{ior}, not
Section~3, and was not touched); the tax-rule shock
$\varepsilon^\tau_t$ (which requires the tax rule's lag/persistence
restructuring described in Subsection~\ref{sec:taxrule}, not made to
the paper's eq.~\ref{tau}); and the COVID labor-supply shock
$\varepsilon^L_t$ (which has no natural home in the paper's derivation,
since it is not tied to any single fixed parameter the way the other
four are).

\begin{center}
\begin{tabular}{@{}lll@{}}
\toprule
Process & Law of motion & Status relative to the paper \\
\midrule
$a_t$ (preference)     & $a_t=a_{t-1}^{\rho_a}\exp(\varepsilon_{a,t})$ & Same functional form; original shock in the paper \\
$z_t$ (productivity)   & $z_t=z_{t-1}^{\rho_z}z^{1-\rho_z}\exp(\varepsilon_{z,t})$ & \textbf{Now in the paper} (eq.~\ref{z}) \\
$v_{a,t}$ (shopping-time pref.) & $v_{a,t}=v_{a,t-1}^{\rho_{va}}v_a^{1-\rho_{va}}\exp(\varepsilon_{va,t})$ & \textbf{Now in the paper} (eq.~\ref{va}) \\
$x_{a,t}$ (deposit tech.) & $x_{a,t}=x_{a,t-1}^{\rho_{xa}}x_a^{1-\rho_{xa}}\exp(\varepsilon_{xa,t})$ & \textbf{Now in the paper} (eq.~\ref{xa}) \\
$g_t$ (govt.\ spending) & $g_t=\rho_g g_{t-1}+(1-\rho_g)g+\varepsilon_{g,t}$ & \textbf{Now in the paper} (eq.~\ref{g}) \\
$z^r_t$ (MP shock)      & $z^r_t=b_r z^r_{t-1}+\varepsilon^r_t$ & Still new: no shock in the paper's rule (\ref{ior}) \\
$\varepsilon^\tau_t$ (tax shock) & iid, enters tax rule directly & Still new: rule (\ref{tau}) is deterministic \\
$\varepsilon^L_t$ (COVID labor supply) & iid, enters labor-leisure FOC & Still new: no counterpart in the paper \\
\bottomrule
\end{tabular}
\end{center}

Six measurement-error shocks (one per non-debt observable) and a
seventh for $bT^{obs}$ are added on top of these -- see Section~4. These
remain entirely absent from the paper, which defines no observable
equations at all (Section~4 below).

\section{Observable/measurement block (entirely new)}

The paper's Sections~1--4.2 define no mapping from model variables to
data at all. The estimated model adds seven measurement equations, each
a log-deviation from steady state plus measurement error:

\begin{center}
\begin{tabular}{@{}lll@{}}
\toprule
Observable & FRED series & Transformation \\
\midrule
$y^{obs}$ & \texttt{GDPC1} / \texttt{OPHNFB} & $\log(\text{GDPC1}/\text{OPHNFB})$, demeaned \\
$d^{obs}$ & \texttt{DPSACBW027SBOG} / \texttt{OPHNFB} & $\log(\text{Deposits}/\text{OPHNFB})$, demeaned \\
$\pi^{obs}$ & \texttt{GDPDEF} & $\log(\text{GDPDEF}_t/\text{GDPDEF}_{t-1})$, demeaned \\
$r^{obs}$ & \texttt{BAA} & $\log(1+\text{BAA}_q/400)$, demeaned \\
$r^{n,obs}$ & \texttt{FEDFUNDS} & $\log((1+\text{FEDFUNDS}/100)^{0.25})$, demeaned \\
$\tau^{obs}$ & \texttt{FGRECPT} / \texttt{GDP} & $\log(\text{Tax receipts}/\text{Nominal GDP})$, demeaned \\
$bT^{obs}$ & \texttt{FYGFDPUN} / \texttt{GDP} & $\log((\text{FYGFDPUN}/1000)/\text{GDP})$, demeaned \\
\bottomrule
\end{tabular}
\end{center}

Two choices are worth flagging as deliberate departures from the most
obvious extension of the model: the loan-rate observable uses Moody's
Baa corporate bond yield rather than the bank prime rate, because Baa
carries genuine credit-spread variation instead of being mechanically
pinned to the federal funds rate plus a fixed spread; and the debt
observable uses federal debt \emph{held by the public}
(\texttt{FYGFDPUN}), not gross/total public debt (\texttt{GFDEBTN}),
divided by nominal GDP (\texttt{GDP}). (An earlier attempt used
\texttt{NGDPDQ} as the GDP denominator; that series does not exist on
FRED at all, and demeaned the debt ratio in levels rather than logs,
inconsistent with the measurement equation above -- both fixed in the
current build.) A currency-in-circulation observable, present in some
earlier drafts of this estimation, is excluded; the model is estimated
on seven observables, not eight.

\section{Calibration changes}

Table~\ref{tab:calibcompare} compares every parameter given a numerical
value in Subsection~4.2 against its estimated-model counterpart.
Frequency matters here (Section~1): the paper's values are not
quarterly.

\begin{center}
\begin{tabular}{@{}lccl@{}}
\toprule
Symbol & Paper (Sec.\ 4.2) & Estimated model & Status \\
\midrule
$\chi$        & 1.0    & 0.5 (recalibrated) & Calibrated in both; value changed \\
$\psi$        & 1.3    & 1.3    & Unchanged \\
$\nu$         & 0.25   & 0.25   & Unchanged \\
$\phi_n$      & 0.0005 & 0.0005 & Unchanged \\
$\phi_p$      & 50     & 50     & Unchanged \\
$\Omega$      & 0.0    & 0.0    & Unchanged \\
$x_n$         & 0.95   & 0.95   & Unchanged \\
$\theta$      & --     & 6.0    & New (goods demand elasticity; not calibrated in Sec.\ 4.2 text) \\
$\beta$       & 0.972  & 0.9967 & Recalibrated for quarterly frequency and the ZLB \\
$v_a$         & 1.138  & 2.5035 (initial) & \textbf{Now estimated}, not calibrated \\
$\eta$        & 1.674  & 1.5294 (initial) & \textbf{Now estimated}, not calibrated \\
$x_a$         & 3.541  & 5.2528 & Still calibrated (only its shock s.d.\ is estimated) \\
$\rho_a$      & 0.9    & 0.95   & Recalibrated; fixed in both (see Section~7 caveat) \\
$g$           & 0.2    & 0.2    & Unchanged \\
$\bar\tau$ (paper: $\tau$) & 0.18 & 0.14 (recalibrated) & Deficit target changed from 2\% to 6\% of GDP \\
$\alpha_i$ (paper), $a_i$ (est.) & $\approx$0.3 (Table~\ref{tab:inf}) & 0.181 (post.; prior mean 0.49) & \textbf{Now estimated}, not calibrated \\
$b_\pi$       & (varied in analysis, no fixed value) & 3.47 (post.; prior mean 1.00) & \textbf{Now estimated}, not calibrated; diffuse prior (see Section~6) \\
$\varphi_b$   & (varied in analysis, threshold $<$0.153) & 0.069 (post.; prior mean 0.05) & \textbf{Now estimated}, not calibrated \\
\bottomrule
\end{tabular}
\end{center}
\label{tab:calibcompare}

Parameters with no counterpart in the paper at all, needed only by the
extensions above: $b_y$ (output-response weight, new), $b_r$
(monetary-policy-shock persistence, new), $\rho_z,\rho_{xa},\rho_{va},\rho_g,\rho_\tau$
(persistence of the newly stochastic processes in Section~3, all
absent as concepts from the paper), $\pi^{ss}$ and $r^{\ell,ss}$ (fixed
nominal anchors used only in the linearized debt equation and the
Taylor-rule normalization, Section~2.3), and $\underline{r}^{\,n}=1$
(the ZLB floor itself).

\section{Parameters moved from calibrated to estimated}

Eight structural parameters are estimated via Bayesian methods rather
than calibrated: $v_a$, $\eta$, $b_\pi$, $b_y$, $a_i$, $b_r$, $\rho_z$,
$\varphi_b$. Of these, only $v_a$, $\eta$, and (implicitly, via the
Section~4.3 stability sweep) $\varphi_b$ appear as fixed numbers
anywhere in the paper's calibrated model; $b_\pi$, $b_y$, $a_i$, $b_r$,
and $\rho_z$ either have no fixed calibrated value in the paper (they
are swept over ranges for the determinacy analysis, or -- $b_y$, $b_r$
-- do not exist in the paper's model at all).

\textbf{Prior revision (current round).} $b_\pi$ and $b_y$'s priors were
substantially loosened -- from Normal(2.10, 0.15) and Gamma(0.28, 0.10)
to Gamma(1.00, 1.00) and Gamma(1.00, 0.70) -- so that the posterior on
the Taylor-rule inflation-response coefficient is not mechanically
pulled toward the earlier round's own tight prior; $\varphi_b$'s prior
mean was also lowered, from 0.10 to 0.05. Priors and current-round
posterior results:

\begin{center}
\small
\begin{tabular}{@{}llrrl@{}}
\toprule
Symbol & Prior & Prior mean/std & Post.\ mean (90\% HPD) & \\
\midrule
$v_a$ & Gamma & 2.50 / 0.75 & 2.890\ [2.055, 3.482] \\
$\eta$ & Gamma & 1.53 / 0.40 & 1.136\ [0.922, 1.326] \\
$b_\pi$ & Gamma & \textbf{1.00 / 1.00} & \textbf{3.474\ [3.344, 3.604]} \\
$b_y$ & Gamma & \textbf{1.00 / 0.70} & 0.083\ [0.057, 0.106] \\
$a_i$ & Beta & 0.49 / 0.10 & 0.181\ [0.169, 0.195] \\
$b_r$ & Beta & 0.36 / 0.10 & 0.262\ [0.171, 0.347] \\
$\rho_z$ & Beta & 0.45 / 0.10 & 0.443\ [0.292, 0.549] \\
$\varphi_b$ & Beta & \textbf{0.05} / 0.05 & 0.069\ [0.030, 0.101] \\
\bottomrule
\end{tabular}
\end{center}

With the informative prior removed, $b_\pi$'s posterior does not fall
toward the diffuse prior's mean of 1 -- it moves to 3.47, \emph{higher}
than the earlier tight-prior estimate of 2.17. The data want a stronger
inflation response than the previous prior allowed; the conclusion that
the Taylor principle holds is a property of the data, not an artifact of
prior choice. $b_y$'s prior needed reshaping once during this process:
an initial Gamma(1.00, 1.00) (shape 1, i.e.\ Exponential, whose density
is maximized at zero) put $b_y$'s posterior mode at exactly 0.0000 with
a non-positive-definite Hessian there -- a corner solution the prior
itself pre-loads. Raising the shape via Gamma(1.00, 0.70) (mode near
0.51, no boundary pre-loading) resolved this without narrowing the
prior's effective diffuseness. $a_i$ is essentially unchanged from the
previous round (0.181 vs.\ 0.180) despite its own prior being untouched,
a useful cross-round consistency check.

$\rho_a$ and $\rho_{xa}$ were also tried as free parameters in an
earlier attempt; freeing them drove $\rho_a$ to $\approx$0.9992 --
effectively a unit root -- and produced a non-positive-definite Hessian
at the posterior mode, so both remain fixed at 0.95 in all reported
results.

\section{Estimation methodology (absent from Sections 1--4.2 by construction)}

Sections~1--3 and Subsections~4.1--4.2 describe only a calibrated,
deterministic-steady-state model; nothing about Bayesian inference
appears until Section~5 (which, as noted in the abstract, is itself
outdated relative to the results actually obtained). Briefly, for a
replacement Section~5 to describe accurately:

\begin{itemize}
\item \textbf{Two-stage mode-finding.} OccBin's piecewise-linear
likelihood is nonsmooth at ZLB regime switches, so the posterior mode
is first located on the linear (non-OccBin) approximation
(\texttt{mode\_compute=4}), then MCMC is restarted at
\texttt{mode\_compute=0} from that mode with the true OccBin likelihood
switched on.
\item \textbf{MCMC jumping covariance.} An initial run using a diagonal
proposal covariance (\texttt{prior\_variance}) produced chain
disagreement traceable to that covariance ignoring the $b_\pi$/$b_y$
correlation. A Hessian-based covariance fixed this under the original
calibration but showed renewed (milder) disagreement after
recalibrating $\bar\tau,\chi$; that round's reported (``gen1'') results
instead used a jumping covariance built from a preliminary round's own
pooled posterior draws -- standard adaptive-MCMC pilot/tune practice,
verified against Dynare's source to get a subtle sign/inversion
convention right (the saved matrix must be a precision, not a
covariance). The current, diffuse-prior round did not need this extra
step: a plain Hessian-based covariance (evaluated at the Stage-1 mode)
was sufficient throughout once the jump scale was tuned down, and
achieved the best chain-mean agreement of any round for three of the
four policy parameters.
\item \textbf{Acceptance-ratio discipline.} Any chain with an
acceptance ratio below 0.05 is treated as a failed run and re-estimated
with an adjusted jump scale; short (3,000-draw) diagnostics
systematically overstated what a full 50,000-draw run would sustain, so
later rounds built in extra margin before committing to a full run. The
current round's first production attempt cleared the 0.05 floor
(5.4\%/5.3\%) but only barely, with weak Geweke convergence -- rather
than accept a technically-passing but marginal result, it was treated as
a failure, and a longer (8,000-draw) diagnostic was used before the next
full run, since the usual 3,000-draw diagnostic had just mispredicted
that round's production behavior especially badly ($\sim$5$\times$
degradation, worse than any earlier round).
\item \textbf{Smoother limitation.} OccBin's nonlinear smoother did not
reach a fixed point in any round attempted, so no OccBin-consistent
smoothed shock/state decomposition is available; this does not affect
the posterior parameter estimates, which are informed by OccBin's
\emph{filter} (used throughout MCMC) rather than the smoother.
\end{itemize}

The full tuning history, Geweke convergence diagnostics, determinacy/
E-stability results at the estimated posterior, and steady-state
implications are documented in
\texttt{mmt8bOccBin3\_covid\_bT\_baa\_epsL\_summary.tex} and are not
reproduced here.

\end{document}
